\addcontentsline{toc}{chapter}{СПИСОК ИСПОЛЬЗОВАННЫХ ИСТОЧНИКОВ}
\begin{thebibliography}{}
	\bibitem{lit1}\textit {РЕШЕНИЕ ЗАДАЧ С ИСПОЛЬЗОВАНИЕМ РЕКУРСИИ} [Электронный ресурс]. URL: https://al.cs.msu.ru/files/recursia.pdf (дата обращения: 26.10.2025).
	\bibitem{lit2}\textit {Рекурсия} [Электронный ресурс]. URL: 
	https://zns.susu.ru/IT/BP/10.pdf(дата обращения: 26.10.2025).
	\bibitem{lit3}\textit {Выполнение программы. Программный стек. Стековые кадры} [Электронный ресурс]. URL: 
	https://it.kgsu.ru/Python\_Recursion/pythonRecursion159.html (дата обращения: 26.10.2025).
	\bibitem{lit4}М. В. Ульянов. РЕСУРСНО-ЭФФЕКТИВНЫЕ КОМПЬЮТЕРНЫЕ АЛГОРИТМЫ. РАЗРАБОТКА И АНАЛИЗ — Наука Физмалит 2007.
	\bibitem{lit5}\textit {Standard C++} [Электронный ресурс]. URL: 
	https://isocpp.org/ (дата обращения: 26.10.2025).
	\bibitem{lit6}\textit {Python} [Электронный ресурс]. URL: https://www.python.org/ (дата обращения: 26.10.2025).
	\bibitem{lit7}\textit {Pandas} [Электронный ресурс]. URL: https://pandas.pydata.org/ (дата обращения: 26.10.2025).
	\bibitem{lit8}\textit {Matplotlib — Visualization with Python} [Электронный ресурс]. URL:  https://matplotlib.org/ (Дата обращения: 26.10.2025).
	\bibitem{lit9}\textit {Clang command line argument reference} [Электронный ресурс]. URL: https://clang.llvm.org/docs/ClangCommandLineReference.html\#optimization-level(Дата обращения: 27.10.2025).
\end{thebibliography}

